\IfFileExists{stacks-project.cls}{%
\documentclass{stacks-project}
}{%
\documentclass{amsart}
}

% For dealing with references we use the comment environment
\usepackage{verbatim}
\newenvironment{reference}{\comment}{\endcomment}
%\newenvironment{reference}{}{}
\newenvironment{slogan}{\comment}{\endcomment}
\newenvironment{history}{\comment}{\endcomment}

% For commutative diagrams we use Xy-pic
\usepackage[all]{xy}

% We use 2cell for 2-commutative diagrams.
\xyoption{2cell}
\UseAllTwocells

% We use multicol for the list of chapters between chapters
\usepackage{multicol}

% This is generally recommended for better output
\usepackage{lmodern}
\usepackage[T1]{fontenc}

% For cross-file-references
\usepackage{xr-hyper}

% Package for hypertext links:
\usepackage{hyperref}

% For any local file, say "hello.tex" you want to link to please
% use \externaldocument[hello-]{hello}
\externaldocument[introduction-]{introduction}
\externaldocument[conventions-]{conventions}
\externaldocument[sets-]{sets}
\externaldocument[categories-]{categories}
\externaldocument[topology-]{topology}
\externaldocument[sheaves-]{sheaves}
\externaldocument[sites-]{sites}
\externaldocument[stacks-]{stacks}
\externaldocument[fields-]{fields}
\externaldocument[algebra-]{algebra}
\externaldocument[brauer-]{brauer}
\externaldocument[homology-]{homology}
\externaldocument[derived-]{derived}
\externaldocument[simplicial-]{simplicial}
\externaldocument[more-algebra-]{more-algebra}
\externaldocument[smoothing-]{smoothing}
\externaldocument[modules-]{modules}
\externaldocument[sites-modules-]{sites-modules}
\externaldocument[injectives-]{injectives}
\externaldocument[cohomology-]{cohomology}
\externaldocument[sites-cohomology-]{sites-cohomology}
\externaldocument[dga-]{dga}
\externaldocument[dpa-]{dpa}
\externaldocument[sdga-]{sdga}
\externaldocument[hypercovering-]{hypercovering}
\externaldocument[schemes-]{schemes}
\externaldocument[constructions-]{constructions}
\externaldocument[properties-]{properties}
\externaldocument[morphisms-]{morphisms}
\externaldocument[coherent-]{coherent}
\externaldocument[divisors-]{divisors}
\externaldocument[limits-]{limits}
\externaldocument[varieties-]{varieties}
\externaldocument[topologies-]{topologies}
\externaldocument[descent-]{descent}
\externaldocument[perfect-]{perfect}
\externaldocument[more-morphisms-]{more-morphisms}
\externaldocument[flat-]{flat}
\externaldocument[groupoids-]{groupoids}
\externaldocument[more-groupoids-]{more-groupoids}
\externaldocument[etale-]{etale}
\externaldocument[chow-]{chow}
\externaldocument[intersection-]{intersection}
\externaldocument[pic-]{pic}
\externaldocument[weil-]{weil}
\externaldocument[adequate-]{adequate}
\externaldocument[dualizing-]{dualizing}
\externaldocument[duality-]{duality}
\externaldocument[discriminant-]{discriminant}
\externaldocument[derham-]{derham}
\externaldocument[local-cohomology-]{local-cohomology}
\externaldocument[algebraization-]{algebraization}
\externaldocument[curves-]{curves}
\externaldocument[resolve-]{resolve}
\externaldocument[models-]{models}
\externaldocument[functors-]{functors}
\externaldocument[equiv-]{equiv}
\externaldocument[pione-]{pione}
\externaldocument[etale-cohomology-]{etale-cohomology}
\externaldocument[proetale-]{proetale}
\externaldocument[relative-cycles-]{relative-cycles}
\externaldocument[more-etale-]{more-etale}
\externaldocument[trace-]{trace}
\externaldocument[crystalline-]{crystalline}
\externaldocument[spaces-]{spaces}
\externaldocument[spaces-properties-]{spaces-properties}
\externaldocument[spaces-morphisms-]{spaces-morphisms}
\externaldocument[decent-spaces-]{decent-spaces}
\externaldocument[spaces-cohomology-]{spaces-cohomology}
\externaldocument[spaces-limits-]{spaces-limits}
\externaldocument[spaces-divisors-]{spaces-divisors}
\externaldocument[spaces-over-fields-]{spaces-over-fields}
\externaldocument[spaces-topologies-]{spaces-topologies}
\externaldocument[spaces-descent-]{spaces-descent}
\externaldocument[spaces-perfect-]{spaces-perfect}
\externaldocument[spaces-more-morphisms-]{spaces-more-morphisms}
\externaldocument[spaces-flat-]{spaces-flat}
\externaldocument[spaces-groupoids-]{spaces-groupoids}
\externaldocument[spaces-more-groupoids-]{spaces-more-groupoids}
\externaldocument[bootstrap-]{bootstrap}
\externaldocument[spaces-pushouts-]{spaces-pushouts}
\externaldocument[spaces-chow-]{spaces-chow}
\externaldocument[groupoids-quotients-]{groupoids-quotients}
\externaldocument[spaces-more-cohomology-]{spaces-more-cohomology}
\externaldocument[spaces-simplicial-]{spaces-simplicial}
\externaldocument[spaces-duality-]{spaces-duality}
\externaldocument[formal-spaces-]{formal-spaces}
\externaldocument[restricted-]{restricted}
\externaldocument[spaces-resolve-]{spaces-resolve}
\externaldocument[formal-defos-]{formal-defos}
\externaldocument[defos-]{defos}
\externaldocument[cotangent-]{cotangent}
\externaldocument[examples-defos-]{examples-defos}
\externaldocument[algebraic-]{algebraic}
\externaldocument[examples-stacks-]{examples-stacks}
\externaldocument[stacks-sheaves-]{stacks-sheaves}
\externaldocument[criteria-]{criteria}
\externaldocument[artin-]{artin}
\externaldocument[quot-]{quot}
\externaldocument[stacks-properties-]{stacks-properties}
\externaldocument[stacks-morphisms-]{stacks-morphisms}
\externaldocument[stacks-limits-]{stacks-limits}
\externaldocument[stacks-cohomology-]{stacks-cohomology}
\externaldocument[stacks-perfect-]{stacks-perfect}
\externaldocument[stacks-introduction-]{stacks-introduction}
\externaldocument[stacks-more-morphisms-]{stacks-more-morphisms}
\externaldocument[stacks-geometry-]{stacks-geometry}
\externaldocument[moduli-]{moduli}
\externaldocument[moduli-curves-]{moduli-curves}
\externaldocument[examples-]{examples}
\externaldocument[exercises-]{exercises}
\externaldocument[guide-]{guide}
\externaldocument[desirables-]{desirables}
\externaldocument[coding-]{coding}
\externaldocument[obsolete-]{obsolete}
\externaldocument[fdl-]{fdl}
\externaldocument[index-]{index}

% Theorem environments.
%
\theoremstyle{plain}
\newtheorem{theorem}[subsection]{Theorem}
\newtheorem{proposition}[subsection]{Proposition}
\newtheorem{lemma}[subsection]{Lemma}

\theoremstyle{definition}
\newtheorem{definition}[subsection]{Definition}
\newtheorem{example}[subsection]{Example}
\newtheorem{exercise}[subsection]{Exercise}
\newtheorem{situation}[subsection]{Situation}

\theoremstyle{remark}
\newtheorem{remark}[subsection]{Remark}
\newtheorem{remarks}[subsection]{Remarks}

\numberwithin{equation}{subsection}

% Macros
%
\def\lim{\mathop{\mathrm{lim}}\nolimits}
\def\colim{\mathop{\mathrm{colim}}\nolimits}
\def\Spec{\mathop{\mathrm{Spec}}}
\def\Hom{\mathop{\mathrm{Hom}}\nolimits}
\def\Ext{\mathop{\mathrm{Ext}}\nolimits}
\def\SheafHom{\mathop{\mathcal{H}\!\mathit{om}}\nolimits}
\def\SheafExt{\mathop{\mathcal{E}\!\mathit{xt}}\nolimits}
\def\Sch{\mathit{Sch}}
\def\Mor{\mathop{\mathrm{Mor}}\nolimits}
\def\Ob{\mathop{\mathrm{Ob}}\nolimits}
\def\Sh{\mathop{\mathit{Sh}}\nolimits}
\def\NL{\mathop{N\!L}\nolimits}
\def\CH{\mathop{\mathrm{CH}}\nolimits}
\def\proetale{{pro\text{-}\acute{e}tale}}
\def\etale{{\acute{e}tale}}
\def\QCoh{\mathit{QCoh}}
\def\Ker{\mathop{\mathrm{Ker}}}
\def\Im{\mathop{\mathrm{Im}}}
\def\Coker{\mathop{\mathrm{Coker}}}
\def\Coim{\mathop{\mathrm{Coim}}}

% Boxtimes
%
\DeclareMathSymbol{\boxtimes}{\mathbin}{AMSa}{"02}

%
% Macros for moduli stacks/spaces
%
\def\QCohstack{\mathcal{QC}\!\mathit{oh}}
\def\Cohstack{\mathcal{C}\!\mathit{oh}}
\def\Spacesstack{\mathcal{S}\!\mathit{paces}}
\def\Quotfunctor{\mathrm{Quot}}
\def\Hilbfunctor{\mathrm{Hilb}}
\def\Curvesstack{\mathcal{C}\!\mathit{urves}}
\def\Polarizedstack{\mathcal{P}\!\mathit{olarized}}
\def\Complexesstack{\mathcal{C}\!\mathit{omplexes}}
% \Pic is the operator that assigns to X its picard group, usage \Pic(X)
% \Picardstack_{X/B} denotes the Picard stack of X over B
% \Picardfunctor_{X/B} denotes the Picard functor of X over B
\def\Pic{\mathop{\mathrm{Pic}}\nolimits}
\def\Picardstack{\mathcal{P}\!\mathit{ic}}
\def\Picardfunctor{\mathrm{Pic}}
\def\Deformationcategory{\mathcal{D}\!\mathit{ef}}
\usepackage{cleveref}
\crefname{originallemma}{original lemma}{original lemmas}
\Crefname{originallemma}{Original lemma}{Original lemmas}
\makeatletter
\newlabel{algebra-lemma-faithfully-flat-universally-injective}{{05CK}{source}{}{sourcepointer-1}{}}
\newlabel{algebra-lemma-faithfully-flat-universally-injective@cref}{{[originallemma][0][]05CK}{[1][0][]source}{}{}{}}
\makeatother
\crefname{originalequation}{original equation}{original equations}
\Crefname{originalequation}{Original equation}{Original equations}
\makeatletter
\newlabel{equation-localize}{{0ANC}{source}{}{sourcepointer-2}{}}
\newlabel{equation-localize@cref}{{[originalequation][0][]0ANC}{[1][0][]source}{}{}{}}
\makeatother
\crefname{originaldefinition}{original definition}{original definitions}
\Crefname{originaldefinition}{Original definition}{Original definitions}
\makeatletter
\newlabel{example-representable-morphism-from-completion}{{0AN6}{source}{}{sourcepointer-3}{}}
\newlabel{example-representable-morphism-from-completion@cref}{{[originaldefinition][0][]0AN6}{[1][0][]source}{}{}{}}
\makeatother
\crefname{originalreference}{original reference}{original references}
\Crefname{originalreference}{Original reference}{Original references}
\makeatletter
\newlabel{item-axiom-1}{{0AND}{source}{}{sourcepointer-4}{}}
\newlabel{item-axiom-1@cref}{{[originalreference][0][]0AND}{[1][0][]source}{}{}{}}
\makeatother
\crefname{originalreference}{original reference}{original references}
\Crefname{originalreference}{Original reference}{Original references}
\makeatletter
\newlabel{item-axiom-2}{{0ANE}{source}{}{sourcepointer-5}{}}
\newlabel{item-axiom-2@cref}{{[originalreference][0][]0ANE}{[1][0][]source}{}{}{}}
\makeatother
\crefname{originalreference}{original reference}{original references}
\Crefname{originalreference}{Original reference}{Original references}
\makeatletter
\newlabel{item-axiom-3}{{0ANF}{source}{}{sourcepointer-6}{}}
\newlabel{item-axiom-3@cref}{{[originalreference][0][]0ANF}{[1][0][]source}{}{}{}}
\makeatother
\crefname{originallemma}{original lemma}{original lemmas}
\Crefname{originallemma}{Original lemma}{Original lemmas}
\makeatletter
\newlabel{lemma-closed}{{0AMS}{source}{}{sourcepointer-7}{}}
\newlabel{lemma-closed@cref}{{[originallemma][0][]0AMS}{[1][0][]source}{}{}{}}
\makeatother
\crefname{originallemma}{original lemma}{original lemmas}
\Crefname{originallemma}{Original lemma}{Original lemmas}
\makeatletter
\newlabel{lemma-representable-affine}{{0AN5}{source}{}{sourcepointer-8}{}}
\newlabel{lemma-representable-affine@cref}{{[originallemma][0][]0AN5}{[1][0][]source}{}{}{}}
\makeatother
\crefname{originallemma}{original lemma}{original lemmas}
\Crefname{originallemma}{Original lemma}{Original lemmas}
\makeatletter
\newlabel{lemma-representable-by-algebraic-spaces-local}{{0AN4}{source}{}{sourcepointer-9}{}}
\newlabel{lemma-representable-by-algebraic-spaces-local@cref}{{[originallemma][0][]0AN4}{[1][0][]source}{}{}{}}
\makeatother
\crefname{originallemma}{original lemma}{original lemmas}
\Crefname{originallemma}{Original lemma}{Original lemmas}
\makeatletter
\newlabel{lemma-topologically-nilpotent}{{0AMW}{source}{}{sourcepointer-10}{}}
\newlabel{lemma-topologically-nilpotent@cref}{{[originallemma][0][]0AMW}{[1][0][]source}{}{}{}}
\makeatother
\crefname{originallemma}{original lemma}{original lemmas}
\Crefname{originallemma}{Original lemma}{Original lemmas}
\makeatletter
\newlabel{lemma-completion-in-sub}{{0ANB}{source}{}{sourcepointer-11}{}}
\newlabel{lemma-completion-in-sub@cref}{{[originallemma][0][]0ANB}{[1][0][]source}{}{}{}}
\makeatother
\crefname{originaldefinition}{original definition}{original definitions}
\Crefname{originaldefinition}{Original definition}{Original definitions}
\makeatletter
\newlabel{situation-local-property}{{0CBA}{source}{}{sourcepointer-12}{}}
\newlabel{situation-local-property@cref}{{[originaldefinition][0][]0CBA}{[1][0][]source}{}{}{}}
\makeatother
\title{Tautness of an arrow of complete weakly admissible rings is a local property}
\author{The Stacks Project Authors}
\date{stacks-project:formal-spaces.tex:local-snapshot-20261012; GNU FDL 1.2 or later}
\begin{document}
\maketitle
\noindent\textit{Editorial note.} All delivered mathematical text is the author's own, copied from the pinned Stacks Project chapter file formal-spaces.tex; no statement or proof step was rewritten, reordered or removed.\par\medskip
\noindent\textit{Editorial note.} The proof opens by observing that the locality axioms have already been obtained from a geometric statement, and then gives a direct algebraic proof of two of them; the final sentence is the author's 'We omit the proof of (axiom-3).' That sentence is reproduced verbatim and recorded as an obligation, and it is to be read together with the geometric argument earlier in the same span, which does cover the third axiom.\par\medskip
\noindent\textit{Editorial note.} The proof displays a commutative diagram of quotients and a three-line chain of equalities computing an intersection; both lie inside the counted span and are reproduced verbatim.\par\medskip
\noindent\textit{Editorial note.} The statement refers to the notions of tautness and to the locality axioms by label; those references are carried as navigation stubs and the definitions are recorded as prerequisites rather than reproduced.\par\medskip
\noindent\textit{Editorial note.} The counted spans carry 12 navigation stubs; no bibliography citation occurs in the spans.\par\medskip
\noindent\textit{Editorial note.} The preamble span points at the chapter's own preamble.tex shipped with the frozen source, used verbatim as staged.\par\medskip
\noindent\textit{Allowed standard prerequisites (editorial).}\begin{itemize}
\item Tautness of an arrow of complete weakly admissible rings, and complete weakly admissible rings themselves. Same chapter, Definition definition-taut, invoked in the statement at line 5469. Fixes the property whose equivalence with the inverse system formulation and whose locality is asserted.
\item The translation of the equivalent conditions into representability of the corresponding morphism, and locality of representability. Same chapter, Lemma lemma-representable-affine at proof lines 5492-5494 and Lemma lemma-representable-by-algebraic-spaces-local at proof lines 5495-5497. Gives all three locality axioms by the geometric route, before the algebraic proof is given.
\item The situation for a local property with its axioms, and the example showing the comparison arrows satisfy the condition. Same chapter, Situation situation-local-property and Equation equation-localize referred to at proof lines 5500-5502, and Example example-representable-morphism-from-completion at proof lines 5502-5504. Fixes the shape of the diagram the algebraic proof works with and discharges the case of the comparison arrows.
\item The closure of a weak ideal of definition, and topological nilpotence of elements of such closures. Same chapter, Lemma lemma-closed at proof lines 5541-5542 and Lemma lemma-topologically-nilpotent at proof lines 5544-5546. Identify the sets produced by the computation as closures of ideals and show they are weak ideals of definition.
\item Faithful flatness implies that an intersection can be tested after base change. Algebra, Lemma algebra-lemma-faithfully-flat-universally-injective, invoked at proof lines 5537-5539. Gives the third equality in the computation for the second axiom.
\end{itemize}\medskip
\ifdefined\hypertarget\hypertarget{sourcepointer-1}{}\fi
\noindent\textit{Source reference.} Original lemma 05CK: sources/stacks-chow/tags.txt, source lines 6923--6923. Pointer outside the counted span, invoked once for context; the load-bearing argument of the counted statement is the author's proof reproduced in full in the proof span.\par\medskip
\ifdefined\hypertarget\hypertarget{sourcepointer-2}{}\fi
\noindent\textit{Source reference.} Original equation 0ANC: sources/stacks-formal-spaces/pinned-original-source.tex, source lines 5063--5063. Pointer outside the counted span, invoked once for context; the load-bearing argument of the counted statement is the author's proof reproduced in full in the proof span.\par\medskip
\ifdefined\hypertarget\hypertarget{sourcepointer-3}{}\fi
\noindent\textit{Source reference.} Original definition 0AN6: sources/stacks-formal-spaces/pinned-original-source.tex, source lines 4393--4393. Pointer outside the counted span, invoked once for context; the load-bearing argument of the counted statement is the author's proof reproduced in full in the proof span.\par\medskip
\ifdefined\hypertarget\hypertarget{sourcepointer-4}{}\fi
\noindent\textit{Source reference.} Original reference 0AND: sources/stacks-formal-spaces/pinned-original-source.tex, source lines 5086--5086. Pointer outside the counted span, invoked once for context; the load-bearing argument of the counted statement is the author's proof reproduced in full in the proof span.\par\medskip
\ifdefined\hypertarget\hypertarget{sourcepointer-5}{}\fi
\noindent\textit{Source reference.} Original reference 0ANE: sources/stacks-formal-spaces/pinned-original-source.tex, source lines 5090--5090. Pointer outside the counted span, invoked once for context; the load-bearing argument of the counted statement is the author's proof reproduced in full in the proof span.\par\medskip
\ifdefined\hypertarget\hypertarget{sourcepointer-6}{}\fi
\noindent\textit{Source reference.} Original reference 0ANF: sources/stacks-formal-spaces/pinned-original-source.tex, source lines 5095--5095. Pointer outside the counted span, invoked once for context; the load-bearing argument of the counted statement is the author's proof reproduced in full in the proof span.\par\medskip
\ifdefined\hypertarget\hypertarget{sourcepointer-7}{}\fi
\noindent\textit{Source reference.} Original lemma 0AMS: sources/stacks-formal-spaces/pinned-original-source.tex, source lines 922--922. Pointer outside the counted span, invoked once for context; the load-bearing argument of the counted statement is the author's proof reproduced in full in the proof span.\par\medskip
\ifdefined\hypertarget\hypertarget{sourcepointer-8}{}\fi
\noindent\textit{Source reference.} Original lemma 0AN5: sources/stacks-formal-spaces/pinned-original-source.tex, source lines 4239--4239. Pointer outside the counted span, invoked once for context; the load-bearing argument of the counted statement is the author's proof reproduced in full in the proof span.\par\medskip
\ifdefined\hypertarget\hypertarget{sourcepointer-9}{}\fi
\noindent\textit{Source reference.} Original lemma 0AN4: sources/stacks-formal-spaces/pinned-original-source.tex, source lines 4094--4094. Pointer outside the counted span, invoked once for context; the load-bearing argument of the counted statement is the author's proof reproduced in full in the proof span.\par\medskip
\ifdefined\hypertarget\hypertarget{sourcepointer-10}{}\fi
\noindent\textit{Source reference.} Original lemma 0AMW: sources/stacks-formal-spaces/pinned-original-source.tex, source lines 1134--1134. Pointer outside the counted span, invoked once for context; the load-bearing argument of the counted statement is the author's proof reproduced in full in the proof span.\par\medskip
\ifdefined\hypertarget\hypertarget{sourcepointer-11}{}\fi
\noindent\textit{Source reference.} Original lemma 0ANB: sources/stacks-formal-spaces/pinned-original-source.tex, source lines 5029--5029. Pointer outside the counted span, invoked once for context; the load-bearing argument of the counted statement is the author's proof reproduced in full in the proof span.\par\medskip
\ifdefined\hypertarget\hypertarget{sourcepointer-12}{}\fi
\noindent\textit{Source reference.} Original definition 0CBA: sources/stacks-formal-spaces/pinned-original-source.tex, source lines 5059--5059. Pointer outside the counted span, invoked only from a delivered context span as a forward reference; the load-bearing argument of the counted statement is the author's proof reproduced in full in the proof span.\par\medskip
\begin{definition}
\label{definition-taut}
Let $\varphi : A \to B$ be a continuous map of linearly topologized rings.
We say $\varphi$ is {\it taut}\footnote{This is nonstandard notation.
The definition generalizes to modules, by saying a linearly topologized
$A$-module $M$ is $A$-taut if for every open ideal $I \subset A$ the closure
of $IM$ in $M$ is open and these closures form a fundamental system of
neighbourhoods of $0$ in $M$.}
if for every open ideal $I \subset A$ the closure of the ideal $\varphi(I)B$
is open and these closures form a fundamental system of open ideals.
\end{definition}

\begin{lemma}
\label{lemma-representable-property-rings}
Let $B \to A$ be an arrow of $\textit{WAdm}^{count}$.
The following are equivalent
\begin{enumerate}
\item[(a)] $B \to A$ is taut (Definition \ref{definition-taut}),
\item[(b)] for $B \supset J_1 \supset J_2 \supset J_3 \supset \ldots$
a fundamental system of weak ideals of definitions there exist
a commutative diagram
$$
\xymatrix{
A \ar[r] & \ldots \ar[r] & A_3 \ar[r] & A_2 \ar[r] & A_1 \\
B \ar[r] \ar[u] & \ldots \ar[r] & B/J_3 \ar[r] \ar[u] &
B/J_2 \ar[r] \ar[u] & B/J_1 \ar[u]
}
$$
such that $A_{n + 1}/J_nA_{n + 1} = A_n$ and $A = \lim A_n$
as topological ring.
\end{enumerate}
Moreover, these equivalent conditions define a local property,
i.e., they satisfy axioms (\ref{item-axiom-1}), (\ref{item-axiom-2}),
(\ref{item-axiom-3}).
\end{lemma}

\begin{proof}
The equivalence of (a) and (b) is immediate. Below we will give an
algebraic proof of the axioms, but it turns out we've already proven
them. Namely, using Lemma \ref{lemma-representable-affine}
the equivalent conditions (a) and (b) translate to saying the
corresponding morphism of affine formal algebraic spaces is representable
by algebraic spaces.
Since this condition is ``\'etale local on the source and target'' by
Lemma \ref{lemma-representable-by-algebraic-spaces-local}
we immediately get axioms (\ref{item-axiom-1}), (\ref{item-axiom-2}), and
(\ref{item-axiom-3}).

\medskip\noindent
Direct algebraic proof of (\ref{item-axiom-1}), (\ref{item-axiom-2}),
(\ref{item-axiom-3}). Let a diagram (\ref{equation-localize}) as in
Situation \ref{situation-local-property} be given.
By Example \ref{example-representable-morphism-from-completion}
the maps $A \to (A')^\wedge$ and $B \to (B')^\wedge$
satisfy (a) and (b).

\medskip\noindent
Assume (a) and (b) hold for $\varphi$. Let $J \subset B$ be a weak ideal
of definition. Then the closure of $JA$, resp.\ $J(B')^\wedge$
is a weak ideal of definition $I \subset A$, resp.\ $J' \subset (B')^\wedge$.
Then the closure of $I(A')^\wedge$ is a weak ideal of definition
$I' \subset (A')^\wedge$. A topological argument shows that $I'$ is also
the closure of $J(A')^\wedge$ and of $J'(A')^\wedge$.
Finally, as $J$ runs over a fundamental system of weak ideals of definition
of $B$ so do the ideals $I$ and $I'$ in $A$ and $(A')^\wedge$.
It follows that (a) holds for $\varphi'$. This proves (\ref{item-axiom-1}).

\medskip\noindent
Assume $A \to A'$ is faithfully flat and that (a) and (b) hold for $\varphi'$.
Let $J \subset B$ be a weak ideal of definition. Using (a) and (b)
for the maps $B \to (B')^\wedge \to (A')^\wedge$ we find that the
closure $I'$ of $J(A')^\wedge$ is a weak ideal of definition.
In particular, $I'$ is open and hence the inverse image of $I'$
in $A$ is open. Now we have (explanation below)
\begin{align*}
A \cap I'
& =
A \cap \bigcap (J(A')^\wedge + \Ker((A')^\wedge \to A'/I_0A')) \\
& =
A \cap \bigcap \Ker((A')^\wedge \to A'/JA' + I_0 A') \\
& = \bigcap (JA + I_0)
\end{align*}
which is the closure of $JA$ by Lemma \ref{lemma-closed}.
The intersections are over weak ideals of definition $I_0 \subset A$.
The first equality because a fundamental system of neighbourhoods of
$0$ in $(A')^\wedge$ are the kernels of the maps $(A')^\wedge \to A'/I_0A'$.
The second equality is trivial. The third equality because $A \to A'$
is faithfully flat, see
Algebra, Lemma \ref{algebra-lemma-faithfully-flat-universally-injective}.
Thus the closure of $JA$ is open. By Lemma \ref{lemma-topologically-nilpotent}
the closure of $JA$
is a weak ideal of definition of $A$. Finally, given a weak
ideal of definition $I \subset A$ we can find $J$ such that
$J(A')^\wedge$ is contained in the closure of $I(A')^\wedge$
by property (a) for $B \to (B')^\wedge$ and $\varphi'$.
Thus we see that (a) holds for $\varphi$. This proves (\ref{item-axiom-2}).

\medskip\noindent
We omit the proof of (\ref{item-axiom-3}).
\end{proof}

\end{document}
